\documentclass{article}
\usepackage{amsmath,amssymb,amsthm}
\usepackage{algorithm}
\usepackage{algpseudocode}
\usepackage{hyperref}
\usepackage{xcolor}
\newtheorem{theorem}{Theorem}
\newtheorem{lemma}{Lemma}
\newtheorem{assumption}{Assumption}
\newtheorem{corollary}{Corollary}
\newcommand{\sign}{\operatorname{sign}}
\newcommand{\E}{\mathbb{E}}
\newcommand{\R}{\mathbb{R}}

\begin{document}

\section*{SignSGD with Median-of-Gradients}

\section*{Mathematical Formulation}
Let $f:\R^{d}\rightarrow\R$ be the (possibly non‑convex) loss function we wish to minimise.
At iteration $k$ we have a current iterate $x_k\in\R^{d}$.
A mini‑batch of stochastic samples is drawn on $M$ parallel workers.
Each worker $m\in\{1,\dots,M\}$ computes a stochastic gradient
\[
g_k^{(m)} \;\in\; \R^{d},\qquad 
\mathbb{E}\!\left[g_k^{(m)}\mid x_k\right]=\nabla f(x_k).
\]
For each coordinate $i\in\{1,\dots,d\}$ we take the \emph{median} of the $M$ scalar
values $\{g_{k,i}^{(1)},\dots,g_{k,i}^{(M)}\}$ and denote it by
\[
\tilde g_{k,i}\;:=\;\operatorname{median}\!\bigl(g_{k,i}^{(1)},\dots,g_{k,i}^{(M)}\bigr).
\]
The update direction is the sign of this median vector:
\[
d_k \;:=\; \sign\!\bigl(\tilde g_k\bigr),\qquad 
\tilde g_k:=\bigl(\tilde g_{k,1},\dots,\tilde g_{k,d}\bigr)^{\!\top}.
\]
With a stepsize (learning rate) $\alpha>0$, the iteration reads
\begin{equation}
\boxed{\,x_{k+1}=x_k-\alpha\,d_k\,=\,x_k-\alpha\,\sign\!\bigl(\tilde g_k\bigr).\,}
\label{eq:update}
\end{equation}
The algorithm is often called **SignSGD with median‑of‑gradients**.

\section*{Pseudocode}
\begin{algorithm}[H]
\caption{SignSGD with Median‑of‑Gradients}
\begin{algorithmic}[1]
\State \textbf{Input:} stepsize $\alpha>0$, number of workers $M$, max iterations $T$.
\State Initialise $x_0\in\R^{d}$.
\For{$k=0,1,\dots,T-1$}
    \For{each worker $m=1,\dots,M$ \textbf{in parallel}}
        \State Sample a mini‑batch $\mathcal{B}_k^{(m)}$.
        \State Compute stochastic gradient $g_k^{(m)}=\nabla f(x_k;\mathcal{B}_k^{(m)})$.
    \EndFor
    \For{$i=1,\dots,d$}
        \State $\tilde g_{k,i}\gets\operatorname{median}\bigl(g_{k,i}^{(1)},\dots,g_{k,i}^{(M)}\bigr)$.
    \EndFor
    \State $d_k\gets \sign(\tilde g_k)$.
    \State $x_{k+1}\gets x_k-\alpha d_k$.
\EndFor
\State \textbf{Output:} $x_T$ (or the average iterate $\frac1T\sum_{k=0}^{T-1}x_k$).
\end{algorithmic}
\end{algorithm}

\section*{Dry Run Example}
Consider the one‑dimensional quadratic
\[
f(w)=(w-3)^2,\qquad \nabla f(w)=2(w-3).
\]
We use a deterministic gradient (so $M=1$), stepsize $\alpha=0.1$ and start at $w_0=0$.
\begin{center}
\begin{tabular}{c|c|c|c}
iteration $k$ & $w_k$ & $\nabla f(w_k)$ & $w_{k+1}=w_k-\alpha\,\sign(\nabla f(w_k))$ \\\hline
0 & $0$ & $2(0-3)=-6$ & $0-0.1\cdot(-1)=0.1$\\
1 & $0.1$ & $2(0.1-3)=-5.8$ & $0.1-0.1\cdot(-1)=0.2$\\
2 & $0.2$ & $2(0.2-3)=-5.6$ & $0.2-0.1\cdot(-1)=0.3$\\
\end{tabular}
\end{center}
The objective values are $f(0)=9$, $f(0.1)=8.41$, $f(0.2)=7.84$, $f(0.3)=7.29$,
showing a monotone decrease despite using only the sign information.

\newpage
\section*{Assumptions}
\begin{assumption}[Smoothness]
\label{ass:smooth}
$f$ is $L$‑smooth, i.e.
\[
\forall x,y\in\R^{d}:\quad 
f(y)\le f(x)+\langle\nabla f(x),y-x\rangle+\frac{L}{2}\|y-x\|_2^{2}.
\]
\end{assumption}
\begin{assumption}[Lower Boundedness]
\label{ass:lower}
There exists $f_{\inf}>-\infty$ such that $f(x)\ge f_{\inf}$ for all $x$.
\end{assumption}
\begin{assumption}[Unbiased Stochastic Gradients]
\label{ass:unbiased}
For every worker $m$,
\[
\E\!\left[g_k^{(m)}\mid x_k\right]=\nabla f(x_k).
\]
\end{assumption}
\begin{assumption}[Median Sign Concentration]
\label{ass:median}
Define for each coordinate $i$ the probability
\[
p_{k,i}:=\Pr\!\bigl(\sign(\tilde g_{k,i})=\sign(\nabla_i f(x_k))\mid x_k\bigr).
\]
Assume there exists a constant $\gamma\in(0,\tfrac12]$ such that
\[
p_{k,i}\;\ge\;\tfrac12+\gamma\qquad\forall k,\;i.
\tag{A4}
\]
Equivalently,
\[
\Pr\!\bigl(\sign(\tilde g_{k,i})\neq\sign(\nabla_i f(x_k))\mid x_k\bigr)\le\tfrac12-\gamma.
\]
\end{assumption}
\begin{assumption}[Bounded Dimensionality]
\label{ass:dim}
The dimension $d$ is finite (no explicit bound needed; it appears in the final rate).
\end{assumption}

\section*{Main Convergence Theorem}
\begin{theorem}[Convergence of SignSGD with Median]
\label{thm:main}
Let Assumptions \ref{ass:smooth}–\ref{ass:dim} hold and use a constant stepsize
$\alpha\in\bigl(0,\frac{2\gamma}{L d}\bigr]$.
Define $\Delta:=f(x_0)-f_{\inf}$.
Then the iterates generated by \eqref{eq:update} satisfy
\[
\frac{1}{T}\sum_{k=0}^{T-1}\E\!\bigl[\|\nabla f(x_k)\|_1\bigr]
\;\le\;
\frac{\Delta}{\alpha\gamma T}
\;+\;
\frac{L\alpha d}{2\gamma}.
\tag{1}
\]
In particular, choosing $\alpha=\sqrt{\frac{2\Delta}{L d\,T}}$ yields
\[
\frac{1}{T}\sum_{k=0}^{T-1}\E\!\bigl[\|\nabla f(x_k)\|_1\bigr]
\;=\;
O\!\bigl(T^{-1/2}\bigr).
\]
\end{theorem}

\section*{Theoretical Properties}
\begin{itemize}
\item \textbf{Stability.} The bound in \eqref{eq:update} contains the term
$\frac{L\alpha d}{2\gamma}$; choosing $\alpha$ too large makes this term dominate,
causing divergence. Hence the stepsize must respect the upper bound in Theorem~\ref{thm:main}.
\item \textbf{Hyperparameter Sensitivity.} The convergence rate depends linearly on
the dimension $d$ (through the $\ell_1$ norm of the gradient) and inversely on the
sign‑correctness margin $\gamma$. Better median concentration (larger $\gamma$)
improves the rate.
\item \textbf{Comparison to SGD.} Classical SGD under the same smoothness
assumptions attains an $O(T^{-1/2})$ bound on $\frac{1}{T}\sum\E\|\nabla f(x_k)\|_2^{2}$.
SignSGD trades the $\ell_2$ norm for the $\ell_1$ norm and reduces communication
to one bit per coordinate, at the price of a factor $d$ in the bound.
\end{itemize}

\section*{Discussion}
The theorem shows that even though only the sign of a (median‑aggregated) stochastic
gradient is used, the algorithm still makes expected progress towards a stationary
point. The median aggregation is crucial: it guarantees the sign‑correctness
probability $p_{k,i}$ is uniformly bounded away from $1/2$, which is precisely the
condition (A4) required for the proof. Without this concentration, the bound would
degenerate to the trivial case where the sign is random.

\newpage
\section*{Preliminaries (Definitions and Lemmas)}
\begin{lemma}[Smoothness Inequality]
\label{lem:smooth}
If $f$ satisfies Assumption~\ref{ass:smooth}, then for any $x$ and any direction $d$,
\[
f(x-\alpha d)\le f(x)-\alpha\langle\nabla f(x),d\rangle+\frac{L\alpha^{2}}{2}\|d\|_{2}^{2}.
\]
\end{lemma}
\begin{proof}
Apply the definition of $L$‑smoothness with $y=x-\alpha d$.
\end{proof}

\begin{lemma}[Sign Expectation under (A4)]
\label{lem:signexp}
For any coordinate $i$,
\[
\E\!\bigl[\sign(\tilde g_{k,i})\mid x_k\bigr]
= (2p_{k,i}-1)\,\sign(\nabla_i f(x_k)).
\]
Consequently, using (A4) we have
\[
\E\!\bigl[\sign(\tilde g_{k,i})\mid x_k\bigr]
\ge 2\gamma\,\sign(\nabla_i f(x_k)).
\tag{2}
\]
\end{lemma}
\begin{proof}
Condition on $x_k$.
With probability $p_{k,i}$ the sign matches $\sign(\nabla_i f(x_k))$, otherwise it is opposite.
Hence
\[
\E[\sign(\tilde g_{k,i})\mid x_k]
= p_{k,i}\cdot\sign(\nabla_i f(x_k)) + (1-p_{k,i})\cdot\bigl(-\sign(\nabla_i f(x_k))\bigr)
=(2p_{k,i}-1)\sign(\nabla_i f(x_k)).
\]
Since $p_{k,i}\ge\tfrac12+\gamma$, we obtain $(2p_{k,i}-1)\ge2\gamma$, proving \eqref{2}.
\end{proof}

\section*{Main Proof (Step‑by‑Step Derivation)}
We prove Theorem~\ref{thm:main}. All expectations are taken w.r.t. the randomness
generated up to iteration $k$ unless otherwise stated.

\noindent\textbf{[Step 1]} Apply Lemma~\ref{lem:smooth} with $d=d_k=\sign(\tilde g_k)$:
\[
f(x_{k+1})\le f(x_k)-\alpha\langle\nabla f(x_k),d_k\rangle
+\frac{L\alpha^{2}}{2}\|d_k\|_{2}^{2}.
\tag{3}
\]

\noindent\textbf{[Step 2]} Since each component of $d_k$ equals $\pm1$,
\[
\|d_k\|_{2}^{2}= \sum_{i=1}^{d} 1 = d.
\tag{4}
\]

\noindent\textbf{[Step 3]} Take conditional expectation $\E[\cdot\mid x_k]$ on \eqref{3}:
\[
\E\!\bigl[f(x_{k+1})\mid x_k\bigr]
\le f(x_k)-\alpha\,\E\!\bigl[\langle\nabla f(x_k),d_k\rangle\mid x_k\bigr]
+\frac{L\alpha^{2}d}{2}.
\tag{5}
\]

\noindent\textbf{[Step 4]} Expand the inner product coordinate‑wise and use Lemma~\ref{lem:signexp}:
\[
\E\!\bigl[\langle\nabla f(x_k),d_k\rangle\mid x_k\bigr]
= \sum_{i=1}^{d}\nabla_i f(x_k)\,
\E\!\bigl[\sign(\tilde g_{k,i})\mid x_k\bigr]
\ge \sum_{i=1}^{d}\nabla_i f(x_k)\,
\bigl(2\gamma\,\sign(\nabla_i f(x_k))\bigr)
=2\gamma\sum_{i=1}^{d}|\nabla_i f(x_k)|
=2\gamma\|\nabla f(x_k)\|_{1}.
\tag{6}
\]

\noindent\textbf{[Step 5]} Substitute \eqref{6} into \eqref{5}:
\[
\E\!\bigl[f(x_{k+1})\mid x_k\bigr]
\le f(x_k)-2\alpha\gamma\|\nabla f(x_k)\|_{1}
+\frac{L\alpha^{2}d}{2}.
\tag{7}
\]

\noindent\textbf{[Step 6]} Take full expectation (remove conditioning) and rearrange:
\[
2\alpha\gamma\,\E\!\bigl[\|\nabla f(x_k)\|_{1}\bigr]
\le \E\!\bigl[f(x_k)-f(x_{k+1})\bigr]
+\frac{L\alpha^{2}d}{2}.
\tag{8}
\]

\noindent\textbf{[Step 7]} Sum \eqref{8} over $k=0,\dots,T-1$:
\[
2\alpha\gamma\sum_{k=0}^{T-1}\E\!\bigl[\|\nabla f(x_k)\|_{1}\bigr]
\le \E\!\bigl[f(x_0)-f(x_T)\bigr]
+\frac{L\alpha^{2}d}{2}\,T.
\tag{9}
\]

\noindent\textbf{[Step 8]} Use the lower bound $f(x_T)\ge f_{\inf}$ and define $\Delta:=f(x_0)-f_{\inf}$:
\[
2\alpha\gamma\sum_{k=0}^{T-1}\E\!\bigl[\|\nabla f(x_k)\|_{1}\bigr]
\le \Delta+\frac{L\alpha^{2}d}{2}\,T.
\tag{10}
\]

\noindent\textbf{[Step 9]} Divide both sides of \eqref{10} by $2\alpha\gamma T$:
\[
\frac{1}{T}\sum_{k=0}^{T-1}\E\!\bigl[\|\nabla f(x_k)\|_{1}\bigr]
\le
\frac{\Delta}{2\alpha\gamma T}
+\frac{L\alpha d}{4\gamma}.
\tag{11}
\]

\noindent\textbf{[Step 10]} Multiply numerator and denominator of the first term by $2$ to obtain the statement of the theorem:
\[
\frac{1}{T}\sum_{k=0}^{T-1}\E\!\bigl[\|\nabla f(x_k)\|_{1}\bigr]
\le
\frac{\Delta}{\alpha\gamma T}
+\frac{L\alpha d}{2\gamma},
\]
which is exactly \eqref{1}.  ∎

\section*{Rate Derivation (With Constants)}
Choose $\alpha$ to balance the two terms on the right‑hand side of \eqref{1}:
\[
\frac{\Delta}{\alpha\gamma T} = \frac{L\alpha d}{2\gamma}
\quad\Longrightarrow\quad
\alpha^{2}= \frac{2\Delta}{L d\,T}
\quad\Longrightarrow\quad
\alpha = \sqrt{\frac{2\Delta}{L d\,T}}.
\]
Plugging this optimal $\alpha$ back into \eqref{1} yields
\[
\frac{1}{T}\sum_{k=0}^{T-1}\E\!\bigl[\|\nabla f(x_k)\|_{1}\bigr]
\le
\sqrt{\frac{2L d\Delta}{\gamma^{2}T}}
\;=\;
O\!\bigl(T^{-1/2}\bigr).
\]

\section*{Special Cases}
\begin{itemize}
\item \textbf{Convex $f$.}
When $f$ is convex, the $\ell_{1}$ bound can be turned into an $\ell_{2}$ bound
via $\|\nabla f\|_{2}\le \|\nabla f\|_{1}$, yielding the same $O(T^{-1/2})$ rate for the
gradient norm.
\item \textbf{Deterministic gradients ($\sigma=0$).}
If $g_k^{(m)}=\nabla f(x_k)$ for all workers, then $p_{k,i}=1$ and $\gamma=1/2$.
The bound simplifies to
\[
\frac{1}{T}\sum_{k=0}^{T-1}\|\nabla f(x_k)\|_{1}
\le
\frac{2\Delta}{\alpha T}+\frac{L\alpha d}{1}.
\]
Choosing $\alpha\propto 1/\sqrt{T}$ still gives $O(T^{-1/2})$.
\item \textbf{Large dimension $d$.}
The linear dependence on $d$ reflects the fact that each iteration moves
$\alpha$ units in every coordinate (sign step).  Using a per‑coordinate stepsize
$\alpha_i=\alpha/d$ would remove the $d$ factor at the expense of slower progress.
\end{itemize}

\end{document}